% ---- Dédicace -------------------------------------------------------
\chapter*{Dédicace}
\addcontentsline{toc}{chapter}{Dédicace}
\thispagestyle{plain}

\vspace*{1.5cm}
\begin{center}
\textcolor{accent}{\rule{0.32\linewidth}{1.2pt}}\\[1.4cm]
\end{center}

\begin{center}\large\itshape
\setstretch{1.6}
À nos chers \textbf{parents},\\
pour leur amour inconditionnel, leurs sacrifices\\
et leur soutien qui n'a jamais faibli.\\[1.1cm]

À nos \textbf{familles} et à nos \textbf{ami(e)s},\\
pour leurs encouragements constants.\\[1.1cm]

À nos \textbf{enseignants},\\
qui ont éveillé en nous la curiosité\\
et le goût du travail bien fait.\\[1.1cm]

À toutes celles et ceux qui croient\\
qu'une donnée, bien analysée, éclaire la décision.
\end{center}

\vspace{1cm}
\begin{center}
\textcolor{accent}{\rule{0.32\linewidth}{1.2pt}}
\end{center}
\begin{flushright}\itshape
\auteurun\ \& \auteurdeux
\end{flushright}

% ---- Remerciements --------------------------------------------------
\chapter*{Remerciements}
\addcontentsline{toc}{chapter}{Remerciements}

\begin{center}\itshape
« La connaissance s'acquiert par l'expérience, et se transmet par l'exemple. »
\end{center}
\vspace{0.8cm}

Au terme de ce projet de fin d'année, nous tenons à exprimer notre profonde
gratitude à toutes les personnes qui ont rendu ce travail possible.

Nos remerciements les plus sincères s'adressent à notre encadrant,
\textbf{\monencadrant}, qui a su transformer ce projet exigeant en une
expérience formatrice remarquable grâce à :
\begin{itemize}
  \item son \textbf{encadrement rigoureux} et ses orientations éclairées ;
  \item sa \textbf{disponibilité} constante et son écoute attentive ;
  \item ses \textbf{conseils avisés}, déterminants pour nos choix techniques ;
  \item son \textbf{exigence bienveillante}, qui nous a poussés à nous dépasser.
\end{itemize}

Son accompagnement a été déterminant pour :
\begin{itemize}
  \item structurer notre démarche d'ingénierie ;
  \item approfondir notre compréhension des architectures multi-agents ;
  \item mener à bien la réalisation pratique de cette plateforme.
\end{itemize}

Nous adressons également notre reconnaissance :
\begin{itemize}
  \item à l'ensemble du \textbf{corps professoral} de l'\monetablissement,
        pour la qualité de la formation dispensée ;
  \item à nos \textbf{familles}, pour leur soutien indéfectible ;
  \item à nos \textbf{camarades de promotion}, pour leur entraide et leurs
        encouragements.
\end{itemize}

Ce projet a constitué une expérience académique enrichissante, permettant une
application concrète des connaissances théoriques. Nous en garderons non
seulement des compétences techniques solides, mais aussi l'exemple d'un
enseignement de grande qualité.

\vspace{0.8cm}
\begin{flushright}\itshape
Avec toute notre reconnaissance,\\[6pt]
\textbf{\auteurun\ \& \auteurdeux}
\end{flushright}

% ---- Résumé ---------------------------------------------------------
\chapter*{Résumé}
\addcontentsline{toc}{chapter}{Résumé}

Les outils de \textit{Business Intelligence} (BI) traditionnels exigent un
travail manuel important : définition des indicateurs, modélisation des
données, conception des tableaux de bord et rédaction des rapports. Ce
processus est long, coûteux et fortement dépendant d'experts. Ce projet
propose \textbf{\montitre}, une plateforme décisionnelle \emph{généraliste}
et \emph{automatisée} capable d'analyser n'importe quel jeu de données
tabulaire (CSV, Excel, JSON) \emph{sans hypothèse préalable} sur les noms de
colonnes ni sur le domaine métier.

La plateforme repose sur une équipe de \textbf{sept agents d'intelligence
artificielle} spécialisés, orchestrés via le \textbf{protocole MCP} (\textit{Model
Context Protocol}). Un serveur MCP expose, sous forme d'outils contrôlés par
permissions, l'ensemble des capacités du système ; chaque agent (Orchestrateur,
Ingestion, Préparation, KPI, Tableau de bord, Publication, Rédaction) n'a accès
qu'aux outils strictement nécessaires à son rôle. Un \textbf{moteur déterministe
adaptatif} profile automatiquement les données, en déduit des indicateurs et des
visualisations pertinents, tandis que les grands modèles de langage (LLM)
curent, priorisent et rédigent l'analyse. Le système produit, par exécution,
un \textbf{tableau de bord interactif}, un \textbf{rapport décisionnel} (page
web et PDF) et un magasin d'artefacts reproductible.

Des mécanismes de \textbf{fiabilité} ont été conçus : filets de sécurité
déterministes garantissant des livrables complets même en cas de défaillance
d'un agent, vérification automatique des chiffres du rapport face aux
indicateurs calculés, et suivi des coûts (jetons, latence). Une étude de cas
sur un jeu de \textbf{17 769 matchs internationaux de football} valide
l'approche : carte mondiale des buts, détection d'anomalies (dont le match
historique Australie 31--0 Samoa américaines) et mise en évidence de
l'avantage du terrain.

\vspace{6pt}
\noindent\textbf{Mots-clés :} Business Intelligence, agents IA, protocole MCP,
grands modèles de langage, automatisation décisionnelle, visualisation de
données, FastAPI, Plotly.

% ---- Abstract -------------------------------------------------------
\chapter*{Abstract}
\addcontentsline{toc}{chapter}{Abstract}

Traditional Business Intelligence (BI) tools require substantial manual
effort: defining indicators, modelling data, designing dashboards and writing
reports. This process is slow, costly and heavily dependent on experts. This
project introduces \textbf{\montitre}, a \emph{generalist} and \emph{automated}
decision-support platform able to analyse any tabular dataset (CSV, Excel,
JSON) \emph{without prior assumptions} about column names or business domain.

The platform relies on a team of \textbf{seven specialised AI agents}
orchestrated through the \textbf{Model Context Protocol (MCP)}. An MCP server
exposes every capability of the system as permission-controlled tools; each
agent (Orchestrator, Ingestion, Data~Prep, KPI, Dashboard, Publishing,
Reporter) can only access the tools strictly required by its role. A
\textbf{deterministic adaptive engine} automatically profiles the data and
derives relevant indicators and visualisations, while Large Language Models
(LLMs) curate, prioritise and write the analysis. Each run yields an
\textbf{interactive dashboard}, a \textbf{BI report} (web page and PDF) and a
reproducible artifact store.

Reliability mechanisms were designed: deterministic safety nets guaranteeing
complete deliverables even when an agent misbehaves, automatic verification of
the report's figures against the computed indicators, and cost tracking
(tokens, latency). A case study on \textbf{17{,}769 international football
matches} validates the approach: a world map of goals, anomaly detection
(including the historic Australia 31--0 American Samoa match) and the
home-advantage effect.

\vspace{6pt}
\noindent\textbf{Keywords:} Business Intelligence, AI agents, MCP protocol,
Large Language Models, decision automation, data visualisation, FastAPI,
Plotly.

\cleardoublepage
